\documentclass[aspectratio=169,10pt]{beamer}
% Shared preamble for the Algorithms and Programming Concepts slides (Beamer, Metropolis)
\usetheme[progressbar=frametitle,numbering=fraction,background=light,block=fill,sectionpage=none]{metropolis}
\usepackage[T1]{fontenc}
\usepackage{FiraSans}
\usepackage[scaled=.85]{FiraMono}
\usepackage{amsmath,amssymb}
\usefonttheme[onlymath]{serif}
\usepackage{booktabs,array,tabularx}
\usepackage{tikz,pgfplots}
\usepackage{xcolor}
\usepackage{listings}
\usepackage{adjustbox}
\usepackage{fontawesome5}
\hyphenpenalty=10000\exhyphenpenalty=10000
\newenvironment{sfig}[1][0.80]{\begin{adjustbox}{max width=\linewidth,max totalheight=#1\textheight,center}}{\end{adjustbox}}
\usetikzlibrary{arrows.meta,positioning,shapes.geometric,shapes.misc,fit,calc,backgrounds,decorations.pathreplacing,patterns,matrix,chains}
\pgfplotsset{compat=1.17}

\definecolor{tdInk}{HTML}{1D2330}
\definecolor{tdMuted}{HTML}{5D6678}
\definecolor{tdLine}{HTML}{D9D5CA}
\definecolor{tdPaper}{HTML}{FBFAF7}
\definecolor{tdBlue}{HTML}{2457D6}
\definecolor{tdGreen}{HTML}{17855A}
\definecolor{tdAmber}{HTML}{B7791F}
\definecolor{tdRed}{HTML}{C8362F}
\definecolor{tdPurple}{HTML}{7A4FD1}
\setbeamercolor{normal text}{fg=tdInk,bg=tdPaper}
\setbeamercolor{alerted text}{fg=tdRed}
\setbeamercolor{example text}{fg=tdGreen}
\setbeamercolor{frametitle}{bg=tdInk,fg=white}
\setbeamercolor{progress bar}{fg=tdBlue,bg=tdLine}
\setbeamercolor{title separator}{fg=tdBlue}
\setbeamercolor{block title}{bg=tdBlue!12,fg=tdInk}
\setbeamercolor{block body}{bg=tdBlue!5,fg=tdInk}
\setbeamercolor{block title alerted}{bg=tdRed!12,fg=tdRed!80!black}
\setbeamercolor{block body alerted}{bg=tdRed!5,fg=tdInk}
\setbeamercolor{block title example}{bg=tdGreen!14,fg=tdGreen!70!black}
\setbeamercolor{block body example}{bg=tdGreen!6,fg=tdInk}
\setbeamerfont{frametitle}{size=\large,series=\bfseries}
\makeatletter\setlength{\metropolis@frametitle@padding}{1.15ex}\makeatother
\setbeamertemplate{navigation symbols}{}
\setbeamerfont{block title}{size=\small}

% ---- code listings (Python 3.14) ----
\lstdefinestyle{py}{language=Python,
  basicstyle=\ttfamily\fontsize{8.4}{10.1}\selectfont\color{tdInk},
  keywordstyle=\color{tdBlue}\bfseries, commentstyle=\color{tdMuted}\itshape, stringstyle=\color{tdGreen},
  morekeywords={match,case,type,self,True,False,None},
  emph={print,len,range,sorted,min,max,sum,enumerate,zip,heappush,heappop,cache},emphstyle=\color{tdPurple},
  showstringspaces=false, columns=fullflexible, keepspaces=true, frame=single, rulecolor=\color{tdLine},
  backgroundcolor=\color{white}, xleftmargin=2pt, xrightmargin=2pt, aboveskip=2pt, belowskip=0pt, framesep=3pt}
\lstdefinestyle{out}{basicstyle=\ttfamily\fontsize{8}{9.6}\selectfont\color{tdInk}, columns=fullflexible, keepspaces=true,
  frame=single, rulecolor=\color{tdGreen!60}, backgroundcolor=\color{tdGreen!6}, xleftmargin=2pt, xrightmargin=2pt, aboveskip=2pt, belowskip=0pt, framesep=3pt}
% \pycode[size]{snippet} shows snippets/<snippet>.py ; \pyout[size]{snippet} shows snippets/<snippet>.out
\newcommand{\pycode}[2][\relax]{{#1\lstinputlisting[style=py]{snippets/#2.py}}}
\newcommand{\pyout}[2][\relax]{{#1\lstinputlisting[style=out]{snippets/#2.out}}}

% ---- tikz styles ----
\tikzset{>={Stealth[length=2.2mm]},
  card/.style={rounded corners=4pt,draw=#1!70!black,fill=#1!10,align=center,inner sep=5pt,font=\small}, card/.default=tdBlue,
  pill/.style={rounded corners=7pt,draw=none,fill=#1,text=white,inner sep=3pt,font=\scriptsize\bfseries,align=center}, pill/.default=tdBlue,
  arr/.style={->,thick,draw=tdMuted},
  cell/.style={draw=tdInk!70,minimum width=7.5mm,minimum height=7.5mm,inner sep=0pt,font=\small\ttfamily,fill=white},
  tn/.style={circle,draw=tdInk!80,fill=white,minimum size=7mm,inner sep=0pt,font=\small\bfseries}}
\newcommand{\key}[1]{\textbf{\color{tdBlue}#1}}
\newcommand{\bad}[1]{\textbf{\color{tdRed}#1}}
\newcommand{\good}[1]{\textbf{\color{tdGreen!70!black}#1}}
\newcommand{\apcfooter}[1]{\setbeamertemplate{footline}{\hfill\tiny\color{tdMuted}Algorithms and Programming Concepts \,|\, #1 \,|\, BUV\hspace{1.6em}\insertframenumber/\inserttotalframenumber\hspace{0.8em}\vspace{4pt}}}
\newcommand{\apctitle}[2]{\title{#1}\subtitle{#2}\author{Dr.\ Quang-Vinh Dang\\ {\small Senior Lecturer (Assoc.\ Prof)}}\institute{British University Vietnam}\date{Algorithms and Programming Concepts \,|\, Year 1 \,|\, Python 3.14}}

\apcfooter{Searching}
\apctitle{Searching}{Linear search, binary search, and why sorted data is powerful}

% ---- drawing helpers; the figures are generated from real traces by searching.py (snippets/tk_*.out)
\tikzset{mcell/.style={cell,minimum width=5.6mm,minimum height=5.6mm,font=\scriptsize\ttfamily}}
% \cellsx{style}{step}{x0}{y}{value/colour,...}
\newcommand{\cellsx}[5]{\foreach \v/\c [count=\i from 0] in {#5}{\node[#1,fill=\c] at (#3+\i*#2,#4) {\v};}}
\newcommand{\rl}[3]{\node[font=\scriptsize,anchor=east,text=tdMuted] at (#1,#2) {#3};}
\newcommand{\rn}[4]{\node[font=\scriptsize,anchor=west,text=#4] at (#1,#2) {#3};}
% \ptr{x}{y}{text}{direction: -1 below, +1 above}{colour}
\newcommand{\ptr}[5]{\node[font=\tiny\bfseries,text=#5] at (#1,#2+0.5*#4) {#3};}
\newcommand{\legsq}[2]{\tikz[baseline=-0.5ex]\node[draw=tdInk!60,fill=#1,minimum size=2.6mm,inner sep=0]{};\;#2}

\begin{document}
\begin{frame}[plain]\titlepage\end{frame}

\begin{frame}{Searching is everywhere}
\begin{columns}[T]
\begin{column}{0.52\textwidth}
{\small
Every time you
\begin{itemize}\setlength\itemsep{2pt}
\item type a name into your phone's contacts,
\item look a student ID up in a database,
\item press Ctrl+F in a document,
\item check whether a username is already taken,
\end{itemize}
a program \key{searches}. Because it happens billions of times a day, \emph{how} it searches decides whether an app feels instant or frozen.}
\end{column}
\begin{column}{0.44\textwidth}
\begin{block}{After this lecture you can}
{\small
\begin{itemize}\setlength\itemsep{2pt}
\item write \textbf{linear search} and \textbf{binary search} in Python;
\item explain \emph{why} binary search needs sorted data;
\item count the steps each one takes;
\item use \texttt{bisect} instead of writing your own;
\item search \textbf{on the answer} (square root, shipping);
\item solve rotated-array, peak and two-pointer problems.
\end{itemize}}
\end{block}
\end{column}
\end{columns}
\vspace{2mm}
{\footnotesize All code in these slides runs on \textbf{Python 3.14}: \texttt{python3.14 searching.py} (the full file is on the course page).}
\end{frame}

\begin{frame}{Linear search: look at every item until you find it}
\begin{columns}[T]
\begin{column}{0.53\textwidth}
\pycode{linear}
\vspace{2mm}
\pyout{linear_demo}
\end{column}
\begin{column}{0.43\textwidth}
{\small
\textbf{Idea.} Start at the first item, compare, move on. Stop when you find the target or run out of items.

\vspace{2mm}
\textbf{Works on any list}: sorted or not. It makes no assumptions.

\vspace{2mm}
\textbf{Returns} the index, or $-1$ as a ``not found'' signal (a common convention, because $-1$ is never a valid index).

\vspace{2mm}
Python's \texttt{x in items} and \texttt{items.index(x)} do exactly this under the hood.}
\end{column}
\end{columns}
\end{frame}

\begin{frame}{How many steps does linear search take?}
\begin{sfig}[0.45]
\begin{tikzpicture}
  \foreach \v [count=\i from 0] in {Mai,Nam,Lan,Binh,Hoa,Tuan,Ha,An}{
    \node[cell,minimum width=11mm,fill=\ifnum\i<2 tdLine!55\else white\fi] at (\i*1.15,0) {\scriptsize\v};
    \node[font=\tiny,text=tdMuted] at (\i*1.15,-0.65) {\i};}
  \draw[->,thick,tdAmber] (2*1.15,1.0) -- (2*1.15,0.5); \node[font=\scriptsize,text=tdAmber!70!black] at (2.3,1.25) {looking at \texttt{Lan}};
  \draw[decorate,decoration={brace,amplitude=4pt,mirror},tdMuted] (-0.5,-0.9) -- (1.65,-0.9) node[midway,below=4pt,font=\tiny] {already checked};
\end{tikzpicture}
\end{sfig}
\vspace{2mm}
\begin{columns}[T]
\column{0.3\textwidth}\centering{\Large\bfseries\color{tdGreen!70!black}1}\\[-1pt]{\footnotesize \textbf{best case}: it is the first item}
\column{0.3\textwidth}\centering{\Large\bfseries\color{tdAmber!70!black}$n/2$}\\[-1pt]{\footnotesize \textbf{average}: about half the list}
\column{0.3\textwidth}\centering{\Large\bfseries\color{tdRed}$n$}\\[-1pt]{\footnotesize \textbf{worst case}: last item, or absent}
\end{columns}
\vspace{2mm}
{\small Double the list, double the work. For a million items, a failed search makes a million comparisons. Can we do better?}
\end{frame}

\begin{frame}{The guessing game: halving the problem}
\begin{columns}[T]
\begin{column}{0.46\textwidth}
{\small
``I am thinking of a number between 1 and 100. After each guess I say \emph{higher} or \emph{lower}.''

\vspace{2mm}
Best strategy: always guess the \textbf{middle}. Each answer throws away half of what is left, so you need at most $\lceil\log_2 100\rceil = 7$ guesses.

\vspace{2mm}
\key{Key requirement:} ``higher/lower'' only helps because the numbers are \emph{in order}. In a shuffled list, a comparison would tell you nothing about where to go next.}
\end{column}
\begin{column}{0.5\textwidth}
\begin{sfig}[0.6]
\begin{tikzpicture}[x=0.075cm]
  \foreach \w/\y/\t in {100/0/{100 candidates},50/-0.9/{50},25/-1.8/{25},13/-2.7/{13},7/-3.6/{7},4/-4.5/{4},2/-5.4/{2},1/-6.3/{1}}{
    \fill[tdBlue!\ifnum\w>60 25\else\ifnum\w>20 40\else\ifnum\w>3 60\else 85\fi\fi\fi] (0,\y) rectangle (\w,\y-0.6);
    \node[font=\scriptsize,anchor=west] at (104,\y-0.3) {\t};}
  \node[font=\scriptsize,text=tdMuted,anchor=west] at (0,-7.2) {after 7 guesses only one number is left};
\end{tikzpicture}
\end{sfig}
\end{column}
\end{columns}
\end{frame}

\begin{frame}{Binary search on a sorted array: target = 47}
\begin{sfig}[0.66]
\begin{tikzpicture}
  \input{snippets/tk_found.out}
\end{tikzpicture}
\end{sfig}
\vspace{1mm}
{\small 16 items, only \textbf{4} comparisons. Grey cells are never looked at again: the window \texttt{[lo, hi]} shrinks by about half each time. (Drawn from the real trace of the code on the next slide.)}
\end{frame}

\begin{frame}{When the target is missing: target = 20}
\begin{sfig}[0.68]
\begin{tikzpicture}
  \input{snippets/tk_missing.out}
\end{tikzpicture}
\end{sfig}
\vspace{1mm}
{\small The window keeps shrinking: after 4 steps \texttt{lo = 5} is larger than \texttt{hi = 4}, the window is \textbf{empty}, so 20 is not in the list. Even a failed search costs only about $\log_2 n$ comparisons.}
\end{frame}

\begin{frame}{Binary search in Python}
\begin{columns}[T]
\begin{column}{0.52\textwidth}
\pycode{binary}
\vspace{2mm}
\pyout{binary_demo}
\end{column}
\begin{column}{0.44\textwidth}
{\small Running an instrumented copy (\texttt{binary\_trace} in the course file) on the array from the previous slides:}
\vspace{1mm}
\pyout{trace_demo}
\vspace{2mm}

{\small Each line is one pass of the loop. \texttt{lo} and \texttt{hi} close in on the target; the loop ends when the middle item is the target (or the window is empty).}
\end{column}
\end{columns}
\end{frame}

\begin{frame}{Getting binary search right}
\begin{columns}[T]
\begin{column}{0.48\textwidth}
\begin{block}{The invariant}
{\small \emph{If the target is in the list at all, it is somewhere in \texttt{items[lo..hi]}.}\\[3pt]
It is true at the start (\texttt{lo=0}, \texttt{hi=n-1}). Each step keeps it true, and the window gets smaller. When the window is empty (\texttt{lo > hi}), the target cannot be there.}
\end{block}
{\small\vspace{2mm}
\textbf{Why \texttt{mid + 1} and \texttt{mid - 1}?} We already looked at \texttt{items[mid]}, so it must leave the window. Writing \texttt{lo = mid} can loop forever.}
\end{column}
\begin{column}{0.48\textwidth}
\begin{alertblock}{Classic bugs}
{\small
\begin{itemize}\setlength\itemsep{2pt}
\item \texttt{while lo < hi}: skips the case of one remaining item.
\item \texttt{lo + hi // 2}: operator precedence! Use \texttt{(lo + hi) // 2}.
\item \texttt{lo = mid} (no \texttt{+1}): infinite loop.
\item Searching an \textbf{unsorted} list: it may return $-1$ for an item that is there.
\end{itemize}}
\end{alertblock}
\end{column}
\end{columns}
\end{frame}

\begin{frame}{Counting the steps: $n$ versus $\log_2 n$}
\begin{columns}[T]
\begin{column}{0.46\textwidth}
\pyout{counting_demo}
\vspace{2mm}
{\small A million items: \bad{1,000,000} comparisons for linear search in the worst case, only \good{20} for binary search. Eight billion items need just 33.}
\vspace{1mm}

{\footnotesize Worst-case steps for binary search $=\lfloor\log_2 n\rfloor+1$, which is \texttt{n.bit\_length()} in Python.}
\end{column}
\begin{column}{0.5\textwidth}
\begin{sfig}[0.7]
\begin{tikzpicture}
\begin{axis}[width=9cm,height=6.2cm,xmode=log,ymode=log,log basis x=2,log basis y=2,xlabel={list size $n$},ylabel={worst-case comparisons},axis lines=left,legend pos=north west,legend style={font=\scriptsize},tick label style={font=\scriptsize},label style={font=\small}]
  \addplot[tdRed,very thick,mark=*,mark size=1.5pt] table[x=n,y=linear]{snippets/counting_data.out};
  \addplot[tdGreen,very thick,mark=square*,mark size=1.5pt] table[x=n,y=binary]{snippets/counting_data.out};
  \legend{linear search,binary search}
\end{axis}
\end{tikzpicture}
\end{sfig}
{\tiny measured by running the code: both axes are logarithmic.}
\end{column}
\end{columns}
\end{frame}

\begin{frame}{Don't write it yourself: the \texttt{bisect} module}
\begin{columns}[T]
\begin{column}{0.53\textwidth}
\pycode{bisect_demo}
\vspace{2mm}
\pyout{bisect_demo}
\end{column}
\begin{column}{0.43\textwidth}
{\small
\texttt{bisect\_left(a, x)} returns the position where \texttt{x} would go to keep \texttt{a} sorted: the index of the first element $\geq x$.

\vspace{2mm}
This gives you more than ``found or not'':
\begin{itemize}\setlength\itemsep{2pt}
\item \textbf{insertion point} for a new value;
\item \textbf{how many} copies: \texttt{bisect\_right - bisect\_left};
\item \textbf{how many are smaller}: \texttt{bisect\_left} itself.
\end{itemize}
It is written in C, so it is also faster than a Python loop.}
\end{column}
\end{columns}
\end{frame}

\begin{frame}{Real timings, one million sorted numbers}
\begin{columns}[T]
\begin{column}{0.46\textwidth}
\pyout{timing}
\vspace{2mm}
{\footnotesize Timings come from running \texttt{searching.py} on the machine that built these slides. Yours will differ, but the \emph{ratio} will not.}
\end{column}
\begin{column}{0.5\textwidth}
{\small
Binary search is thousands of times faster here, but it needs \emph{sorted} data, and sorting costs time too.

\vspace{2mm}
\begin{block}{Sort first?}
{\small Sort once, then every query is cheap. If you search \textbf{many times}, sorting wins quickly (see the break-even above). If you search \textbf{once}, just scan.}
\end{block}
\vspace{1mm}
\textbf{Data changes often?} Keeping a list sorted costs work on every insert. A \texttt{set} or \texttt{dict} (hashing) answers ``is it there?'' in about constant time, no order needed.}
\end{column}
\end{columns}
\end{frame}

\begin{frame}{Choosing a search}
\renewcommand{\arraystretch}{1.25}
\small
\begin{tabularx}{\textwidth}{@{}>{\raggedright\bfseries}p{42mm}XX@{}}
\toprule
Situation & Use & Cost\\
\midrule
Small list, or searching once & \texttt{x in items} (linear search) & $O(n)$\\
Unsorted data, one-off & linear search, or sort then search if you will query again & $O(n)$\\
Sorted data, many queries & binary search / \texttt{bisect} & $O(\log n)$ per query\\
Need ``is it there?'' fast, order irrelevant & \texttt{set} or \texttt{dict} & about $O(1)$\\
Need insertion point, counts, ranges & \texttt{bisect\_left}, \texttt{bisect\_right} & $O(\log n)$\\
\bottomrule
\end{tabularx}
\vspace{3mm}

{\small $O(\cdot)$ is Big-O notation: a way to say how the cost \emph{grows} with $n$. We make it precise in the Big-O lecture.}
\end{frame}

% ===================== worked problems =====================
\begin{frame}{Problem 1: first and last occurrence}
\begin{sfig}[0.62]
\begin{tikzpicture}
  \input{snippets/tk_lower.out}
\end{tikzpicture}
\end{sfig}
\vspace{1mm}
{\small List \texttt{[1,2,2,2,2,3,5,5,8]}, target 2. Plain binary search may stop at \emph{any} of the four 2s. \texttt{lower\_bound} never stops early: the window is \emph{half-open} (blue line = \texttt{lo}, red line = \texttt{hi}, the cell at \texttt{hi} is outside) and the loop ends when they meet at the first index with value $\ge$ target.}
\end{frame}

\begin{frame}{First and last occurrence: code}
\begin{columns}[T]
\begin{column}{0.54\textwidth}
\pycode{lower}
\vspace{1mm}
\pycode{first_last}
\end{column}
\begin{column}{0.42\textwidth}
\pyout{firstlast_demo}
\vspace{2mm}
{\small
\textbf{Two shapes of binary search.}
\begin{itemize}\setlength\itemsep{2pt}
\item Closed window \texttt{[lo, hi]}, \texttt{while lo <= hi}, \texttt{hi = mid - 1}: look for \emph{one exact} match.
\item Half-open \texttt{[lo, hi)}, \texttt{while lo < hi}, \texttt{hi = mid}: look for a \emph{boundary}.
\end{itemize}
The last occurrence is ``where the next larger value starts'', minus one. \texttt{bisect\_left} and \texttt{bisect\_right} do the same (last line).}
\end{column}
\end{columns}
\end{frame}

\begin{frame}{Binary search on the answer}
\begin{columns}[T]
\begin{column}{0.48\textwidth}
\begin{sfig}[0.3]
\begin{tikzpicture}
  \cellsx{cell}{0.8}{0}{0}{F/tdRed!22,F/tdRed!22,F/tdRed!22,F/tdRed!22,F/tdRed!22,T/tdGreen!30,T/tdGreen!30,T/tdGreen!30,T/tdGreen!30,T/tdGreen!30}
  \foreach \i in {0,...,9}{\node[font=\tiny,text=tdMuted] at (\i*0.8,-0.6) {\i};}
  \draw[->,thick,tdAmber] (5*0.8,0.95) -- (5*0.8,0.5); \node[font=\scriptsize,text=tdAmber!70!black] at (5*0.8,1.2) {first true};
\end{tikzpicture}
\end{sfig}
\vspace{2mm}
{\small
Sometimes there is \textbf{no array} to search, only a yes/no question about a number $x$ that is \emph{monotone}: false for small $x$, true from some point on.

\vspace{2mm}
\key{Recipe}
\begin{enumerate}\setlength\itemsep{2pt}
\item Write \texttt{ok(x)}: ``is $x$ big enough?''
\item Choose \texttt{[lo, hi]} that surely contains the answer, with \texttt{ok(hi)} true.
\item \texttt{first\_true} finds the boundary with $O(\log(\text{range}))$ calls to \texttt{ok}.
\end{enumerate}}
\end{column}
\begin{column}{0.5\textwidth}
\pycode{first_true}
\vspace{2mm}
{\small
This is the half-open pattern again: \texttt{hi = mid}, because \texttt{mid} itself may be the answer.

\vspace{2mm}
The next slides use it for a square root and for a scheduling problem.}
\end{column}
\end{columns}
\end{frame}

\begin{frame}{Example: integer square root}
\begin{columns}[T]
\begin{column}{0.52\textwidth}
\begin{sfig}[0.78]
\begin{tikzpicture}
  \input{snippets/tk_isqrt.out}
\end{tikzpicture}
\end{sfig}
\vspace{1mm}
{\small $n=50$: the question is ``is $r^2>50$?''. Red dot: yes, the answer is to the left; green: no, to the right. The dashed line is $\sqrt{50}+1\approx 8.07$, where the answer flips.}
\end{column}
\begin{column}{0.44\textwidth}
\pycode{isqrt}
\vspace{2mm}
\pyout{isqrt_demo}
\vspace{1mm}
{\small
Largest $r$ with $r^2\le n$ equals (first $r$ with $r^2>n$) $-1$. Even for $n\approx10^{12}$ only about 40 squarings are needed, with no floating-point error.}
\end{column}
\end{columns}
\end{frame}

\begin{frame}{Example: ship packages within $D$ days}
\begin{columns}[T]
\begin{column}{0.54\textwidth}
\pycode{ship}
\end{column}
\begin{column}{0.42\textwidth}
{\small Weights $1,2,\dots,10$ are shipped \emph{in this order}; the ship carries up to \texttt{cap} per day. Smallest \texttt{cap} that finishes in 5 days?}
\vspace{1mm}
\pyout{ship_demo}
\vspace{1mm}
{\small
\texttt{days\_needed} is a greedy $O(n)$ scan. More capacity never needs \emph{more} days: the test is monotone in \texttt{cap}.}
\end{column}
\end{columns}
\end{frame}

\begin{frame}{Why it works: days needed against capacity}
\begin{columns}[T]
\begin{column}{0.58\textwidth}
\begin{sfig}[0.8]
\begin{tikzpicture}
\begin{axis}[width=10.5cm,height=6.4cm,xlabel={capacity},ylabel={days needed},axis lines=left,const plot,ymin=0,ymax=8,xmin=8,xmax=56,ytick={1,...,7},legend pos=north east,legend style={font=\scriptsize},tick label style={font=\scriptsize},label style={font=\small}]
  \addplot[tdBlue,very thick] table[x=cap,y=days]{snippets/ship_data.out};
  \addplot[tdRed,dashed,no markers,domain=8:56,samples=2] {5};
  \addplot[only marks,mark=*,mark size=3pt,tdAmber] coordinates {(15,5)};
  \legend{days needed,limit: 5 days,answer: 15}
\end{axis}
\end{tikzpicture}
\end{sfig}
\end{column}
\begin{column}{0.4\textwidth}
{\small
The measured curve from \texttt{days\_needed} is a \textbf{staircase that only goes down}. That is exactly the monotone shape binary search needs: ``too slow'' for small capacities, ``fast enough'' from 15 on.

\vspace{2mm}
The search in the previous slide tried capacities 32, 21, 15, 12, 14: five \texttt{days\_needed} calls instead of scanning all 46 candidates.

\vspace{2mm}
\emph{Same pattern, other problems:} the smallest speed to eat all bananas in $H$ hours, the largest minimum distance between $k$ cows in stalls, the smallest time for $m$ workers to finish a job.}
\end{column}
\end{columns}
\end{frame}

\begin{frame}{Problem 3: search in a rotated sorted array}
\begin{columns}[T]
\begin{column}{0.66\textwidth}
\begin{sfig}[0.7]
\begin{tikzpicture}
  \input{snippets/tk_rotated.out}
\end{tikzpicture}
\end{sfig}
\vspace{1mm}
{\small Array \texttt{[15,18,22,27,3,5,9,12]}: a sorted list rotated by 4. Target 9. Blue = the half we know is sorted.}
\end{column}
\begin{column}{0.32\textwidth}
{\small
The list is not sorted, but \textbf{at least one half of any window is sorted}. At each step:
\begin{enumerate}\setlength\itemsep{2pt}
\item find the sorted half: if \texttt{a[lo] <= a[mid]} it is the left one;
\item test whether the target lies inside that half's range;
\item if yes, keep that half; else keep the other.
\end{enumerate}
Still $O(\log n)$.}
\end{column}
\end{columns}
\end{frame}

\begin{frame}{Rotated array: code}
\begin{columns}[T]
\begin{column}{0.54\textwidth}
\pycode{rotated}
\end{column}
\begin{column}{0.42\textwidth}
\pyout{rotated_demo}
\vspace{2mm}
{\small
Compare with \texttt{binary\_search}: only the \emph{decision} is more careful. The chained comparison \texttt{a[lo] <= target < a[mid]} says ``target is inside the sorted left half'' in one line.

\vspace{2mm}
The last line tests the function on every rotation of every list of size 1 to 11 against the slow \texttt{list.index}: always test tricky index code like this.}
\end{column}
\end{columns}
\end{frame}

\begin{frame}{Problem 4: find a peak element}
\begin{sfig}[0.7]
\begin{tikzpicture}
  \input{snippets/tk_peak.out}
\end{tikzpicture}
\end{sfig}
\vspace{1mm}
{\small A \textbf{peak} is an item not smaller than its neighbours. The list is \emph{not} sorted, yet binary search works: compare \texttt{a[mid]} with \texttt{a[mid+1]}. If the list goes \textbf{uphill} to the right, a peak must exist on the right (the list cannot climb forever); otherwise a peak exists at \texttt{mid} or to its left.}
\end{frame}

\begin{frame}{Peak element: code}
\begin{columns}[T]
\begin{column}{0.52\textwidth}
\pycode{peak}
\vspace{2mm}
\pyout{peak_demo}
\end{column}
\begin{column}{0.44\textwidth}
{\small
\textbf{Why correct?} Invariant: \emph{a peak lies in} \texttt{[lo, hi]}. If \texttt{a[mid] < a[mid+1]}, then \texttt{mid} is not a peak, and following the climb must reach one (or the end of the list, which counts as a peak).

\vspace{2mm}
This is another \emph{half-open style} loop: \texttt{hi = mid}, never \texttt{hi = mid - 1}, because \texttt{mid} itself may be the peak.

\vspace{2mm}
Binary search needs less than sortedness: it needs a question whose answer \textbf{tells you which half to discard}.}
\end{column}
\end{columns}
\end{frame}

\begin{frame}{Problem 5: two numbers that add up to a target (two pointers)}
\begin{columns}[T]
\begin{column}{0.64\textwidth}
\begin{sfig}[0.7]
\begin{tikzpicture}
  \input{snippets/tk_twosum.out}
\end{tikzpicture}
\end{sfig}
\end{column}
\begin{column}{0.34\textwidth}
{\small
Sorted list \texttt{[1,3,4,6,8,11,15]}, target 14. Checking all pairs is $O(n^2)$.

\vspace{2mm}
\textbf{Two pointers:} \texttt{i} at the smallest, \texttt{j} at the largest. If the sum is too small, only a larger left number can help: move \texttt{i} right. If too big, move \texttt{j} left. Every step throws away one item, so $O(n)$.

\vspace{2mm}
Same ``discard what cannot work'' idea as binary search, with a different kind of discard.}
\end{column}
\end{columns}
\end{frame}

\begin{frame}{Two pointers: code}
\begin{columns}[T]
\begin{column}{0.54\textwidth}
\pycode{twosum}
\vspace{1mm}
\pycode{describe}
\end{column}
\begin{column}{0.42\textwidth}
\pyout{twosum_demo}
\vspace{2mm}
{\small
\texttt{describe} uses \texttt{match}: the pattern \texttt{(i, j)} unpacks a pair, \texttt{None} means ``nothing found''.

\vspace{2mm}
Why is it safe to discard? If \texttt{a[i] + a[j] < target}, then \texttt{a[i]} plus anything $\le$ \texttt{a[j]} is smaller still, so \texttt{a[i]} cannot be in any answer with the current or a smaller \texttt{j}.

\vspace{2mm}
A \texttt{set} gives $O(n)$ even for unsorted input.}
\end{column}
\end{columns}
\end{frame}

\begin{frame}{Recap and practice}
\begin{columns}[T]
\begin{column}{0.48\textwidth}
\begin{block}{Take-aways}
{\small
\begin{enumerate}\setlength\itemsep{3pt}
\item \key{Linear search}: any list, $n$ steps in the worst case.
\item \key{Binary search}: sorted list, about $\log_2 n$ steps, because every comparison halves the window.
\item The \textbf{invariant} (``the target is in \texttt{[lo, hi]}'') is how you get the loop right.
\item Use \texttt{bisect} in real code; \texttt{lower\_bound} finds boundaries.
\item \key{Search on the answer}: any monotone yes/no question can be halved.
\item Binary search only needs a test that says which half to discard (rotated array, peak).
\end{enumerate}}
\end{block}
\end{column}
\begin{column}{0.48\textwidth}
\begin{block}{Practice}
{\small
\begin{enumerate}\setlength\itemsep{3pt}
\item Trace \texttt{binary\_search} by hand on \texttt{[2,4,6,8,10,12]} for target \texttt{5}. What are \texttt{lo, hi, mid} at each step?
\item How many comparisons does binary search need for one billion items? Check with \texttt{n.bit\_length()}.
\item Use \texttt{bisect} to find how many scores in a list are below 50.
\item Use \texttt{first\_true} to find $\lceil\sqrt n\rceil$ without \texttt{math}.
\item Run \texttt{search\_rotated} by hand for target 3 on \texttt{[15,18,22,27,3,5,9,12]}.
\end{enumerate}}
\end{block}
\end{column}
\end{columns}
\end{frame}

\begin{frame}{Challenge problems}
\begin{enumerate}\setlength\itemsep{3pt}\small
\item \textbf{Smallest eating speed}: Koko has piles \texttt{[3,6,7,11]} and $H=8$ hours. Each hour she eats up to $k$ bananas from one pile. Find the smallest $k$. (Search on the answer.)
\item \textbf{Find the minimum} of a rotated sorted array without duplicates in $O(\log n)$.
\item \textbf{Rotated array with duplicates}: what goes wrong when \texttt{a[lo] == a[mid] == a[hi]}? What is the worst case now?
\item \textbf{Median of two sorted arrays} in $O(\log(\min(m,n)))$ (hard: binary search on a partition).
\item \textbf{Aggressive cows}: place $k$ cows in sorted stalls to maximise the \emph{smallest} distance between any two cows.
\item \textbf{Three sum}: all triples in a sorted list that add up to 0. Combine a loop with two pointers: $O(n^2)$.
\item \textbf{Closest pair}: in a sorted list find two items whose sum is \emph{closest} to a target.
\item \textbf{Square root with decimals}: use binary search on real numbers to get $\sqrt 2$ to 8 digits. When do you stop?
\end{enumerate}
\end{frame}
\end{document}
